\chapter{Phân Tích \& Thiết Kế Hệ Thống}

\begin{muctieu}
Chương này phân tích các yêu cầu chức năng và phi chức năng, xây dựng sơ đồ ca sử dụng (Use-Case), thiết kế kiến trúc hệ thống 3 tầng kết hợp Cloud/AI, mô tả chi tiết lược đồ cơ sở dữ liệu MongoDB và trình bày giải thuật lập lịch không trùng lặp (Visited Set Graph).
\end{muctieu}

\section{Phân tích Yêu cầu Hệ thống}

\subsection{Yêu cầu Chức năng (\en{Functional Requirements})}
\begin{itemize}
    \item \textbf{F1 - Khám phá Địa phương (\en{Local Exploration}):} Cho phép du khách duyệt xem danh sách 9 tỉnh thành Miền Trung, xem hình ảnh bìa thực tế, tra cứu địa chỉ, phân loại danh mục (\en{attraction}, \en{restaurant}, \en{cafe}, \en{hotel}), đánh giá sao và mở Google Maps chỉ đường.
    \item \textbf{F2 - Tạo Lịch trình Thông minh (\en{AI Trip Planning}):} Thu nhận các tham số đầu vào (Điểm đến, số ngày từ 1 đến 7, ngân sách, số lượng người, địa điểm ưu tiên) và gửi đến AI Engine để tạo lịch trình chi tiết theo từng khung giờ.
    \item \textbf{F3 - Đổi lịch Tránh mưa Thông minh (\en{Weather Re-scheduler}):} Nhận diện khi gặp thời tiết xấu hoặc theo yêu cầu du khách, tự động hoán đổi các hoạt động ngoài trời sang các địa điểm tham quan/ẩm thực trong nhà.
    \item \textbf{F4 - Chia tiền Nhóm (\en{Group Bill Splitter}):} Cung cấp bảng kê khai chi phí nhiều mục (phòng nghỉ, ăn uống, vé vào cổng, phương tiện) và tự động tính toán số tiền mỗi thành viên cần đóng.
    \item \textbf{F5 - Thẻ vé Hành trình Ngoại tuyến (\en{Offline Travel Pass}):} Kết xuất thông tin chuyến đi thành dạng thẻ vé điện tử (\en{Boarding Pass}) có mã QR mô phỏng để du khách chụp ảnh màn hình lưu trữ sử dụng khi không có sóng 4G/Wifi.
    \item \textbf{F6 - Trợ lý Ảo Tư vấn Du lịch (\en{AI Chat Assistant}):} Khung chat trực tuyến hỗ trợ giải đáp tức thì các thắc mắc về kinh nghiệm du lịch, đặc sản và văn hóa Miền Trung.
    \item \textbf{F7 - Quản lý Tài khoản \& Chuyến đi (\en{Account \& My Trips}):} Đăng ký, đăng nhập bảo mật qua JWT, lưu trữ và xem lại các chuyến đi đã tạo trên mọi thiết bị.
\end{itemize}

\subsection{Yêu cầu Phi Chức năng (\en{Non-Functional Requirements})}
\begin{itemize}
    \item \textbf{N1 - Thời gian phản hồi (\en{Performance}):} Tốc độ tải trang ban đầu $< 1$ giây; thời gian AI sinh lịch trình hoàn chỉnh $< 2$ giây.
    \item \textbf{N2 - Tính khả dụng (\en{Usability}):} Giao diện App-First trực quan, hỗ trợ chế độ giả lập điện thoại (\en{Phone Simulator Mode}) giúp thao tác một tay dễ dàng trên màn hình cảm ứng hoặc thuyết trình trên máy tính.
    \item \textbf{N3 - Tính toàn vẹn dữ liệu (\en{Data Integrity}):} 100\% các địa điểm du lịch có địa chỉ thực tế và ảnh chụp sắc nét, không có dữ liệu giữ chỗ (\en{placeholder}) rỗng.
\end{itemize}

\section{Sơ đồ Ca Sử Dụng (Use-Case Diagram)}

\begin{figure}[H]
\centering
\resizebox{0.95\textwidth}{!}{
\begin{tikzpicture}[
    actor/.style={rectangle, draw=blue!80, fill=blue!15, rounded corners=2mm, thick, minimum width=2.6cm, minimum height=1.1cm, align=center, font=\bfseries},
    usecase/.style={ellipse, draw=orange!80!black, fill=orange!10, thick, minimum width=4.5cm, minimum height=0.9cm, align=center, font=\footnotesize},
    line/.style={draw, thick, -{Latex[length=2mm]}}
]

% Actors
\node [actor] (user) at (0, 0) {Khách Du Lịch\\(User)};
\node [actor] (ai_actor) at (11, 0) {AI Engine\\(Gemini API)};

% Use cases
\node [usecase] (uc1) at (5.5, 3.2) {Khám phá địa danh 9 tỉnh Miền Trung};
\node [usecase] (uc2) at (5.5, 1.9) {Lập lịch trình du lịch thông minh};
\node [usecase] (uc3) at (5.5, 0.6) {Đổi lịch trình tránh mưa tự động};
\node [usecase] (uc4) at (5.5, -0.7) {Chia tiền nhóm (Bill Splitter)};
\node [usecase] (uc5) at (5.5, -2.0) {Xuất thẻ vé hành trình Offline};
\node [usecase] (uc6) at (5.5, -3.3) {Đăng nhập / Lưu chuyến đi};

% User associations
\draw [line] (user) -- (uc1);
\draw [line] (user) -- (uc2);
\draw [line] (user) -- (uc3);
\draw [line] (user) -- (uc4);
\draw [line] (user) -- (uc5);
\draw [line] (user) -- (uc6);

% AI associations
\draw [line] (ai_actor) -- (uc2);
\draw [line] (ai_actor) -- (uc3);

\end{tikzpicture}
}
\caption{Sơ đồ Ca Sử Dụng (Use-Case Diagram) của hệ thống WanderWise}
\label{fig:usecase_detail}
\end{figure}

\section{Kiến trúc Tổng thể Hệ thống (3-Tier Architecture)}

\begin{figure}[H]
\centering
\resizebox{\textwidth}{!}{
\begin{tikzpicture}[
    box/.style={rectangle, draw=#1!80!black, fill=#1!12, rounded corners=3mm, thick, align=center, minimum height=1.3cm, font=\small},
    arrow/.style={draw, -{Latex[length=3mm]}, thick}
]

% Tier 1: Client
\node [box=blue, text width=4.0cm] (client) at (0, 0) {
    \textbf{TẦNG CLIENT}\\[0.1cm]
    Vue.js 3 + Vite SPA\\
    \footnotesize (Mobile \& PC Desktop)
};

% Tier 2: Backend
\node [box=green!70!black, text width=4.4cm] (backend) at (5.8, 0) {
    \textbf{TẦNG BACKEND}\\[0.1cm]
    Node.js + Express API\\
    \footnotesize JWT Auth + Visited Set Graph
};

% Tier 3: Cloud & AI
\node [box=purple, text width=3.8cm] (ai) at (11.5, 2.2) {
    \textbf{AI ENGINE}\\[0.1cm]
    Google Gemini 2.5 Flash
};

\node [box=purple, text width=3.8cm] (db) at (11.5, 0) {
    \textbf{DATABASE CLOUD}\\[0.1cm]
    MongoDB Atlas Cluster
};

\node [box=purple, text width=3.8cm] (weather) at (11.5, -2.2) {
    \textbf{WEATHER API}\\[0.1cm]
    Open-Meteo Realtime
};

% Connections
\draw [arrow] (client) -- node[above]{\footnotesize REST API} node[below]{\footnotesize (JSON)} (backend);
\draw [arrow] (backend) -- node[above, sloped]{\footnotesize Prompt} (ai);
\draw [arrow] (backend) -- node[above]{\footnotesize Mongoose} (db);
\draw [arrow] (backend) -- node[below, sloped]{\footnotesize Coordinates} (weather);

\end{tikzpicture}
}
\caption{Sơ đồ Kiến trúc Tổng thể Hệ thống WanderWise (3-Tier Architecture)}
\label{fig:architecture_detail}
\end{figure}

\section{Thiết kế Cơ sở Dữ liệu MongoDB}

\begin{table}[H]
\centering
\caption{Cấu trúc thực thể Place (Địa điểm du lịch \& Ẩm thực)}
\begin{tabular}{|l|l|c|p{7.5cm}|}
\hline
\textbf{Trường dữ liệu} & \textbf{Kiểu dữ liệu} & \textbf{Bắt buộc} & \textbf{Ý nghĩa} \\ \hline
\code{\_id} & ObjectId & Có & Khóa chính định danh duy nhất \\ \hline
\code{name} & String & Có & Tên địa danh (ví dụ: Bà Nà Hills, Động Phong Nha) \\ \hline
\code{destination} & String & Có & Tỉnh/Thành phố (Đà Nẵng, Quảng Bình, Huế,...) \\ \hline
\code{type} & String & Có & Phân loại: \code{attraction}, \code{restaurant}, \code{cafe}, \code{hotel} \\ \hline
\code{address} & String & Có & Địa chỉ cụ thể số nhà, tên đường, phường, quận \\ \hline
\code{image} & String & Có & Đường dẫn URL ảnh thực tế chất lượng cao \\ \hline
\code{rating} & Number & Có & Điểm đánh giá chất lượng (1.0 đến 5.0) \\ \hline
\code{estimated\_cost} & Number & Có & Chi phí tham quan / suất ăn ước tính (VNĐ) \\ \hline
\code{tags} & Array[String] & Không & Tập từ khóa đặc trưng (ví dụ: bãi biển, di sản, check-in) \\ \hline
\code{latitude} & Number & Có & Tọa độ vĩ độ địa lý \\ \hline
\code{longitude} & Number & Có & Tọa độ kinh độ địa lý \\ \hline
\end{tabular}
\end{table}

\begin{table}[H]
\centering
\caption{Cấu trúc thực thể Trip (Lịch trình chuyến đi)}
\begin{tabular}{|l|l|c|p{7.5cm}|}
\hline
\textbf{Trường dữ liệu} & \textbf{Kiểu dữ liệu} & \textbf{Bắt buộc} & \textbf{Ý nghĩa} \\ \hline
\code{\_id} & ObjectId & Có & Khóa chính \\ \hline
\code{userId} & ObjectId & Có & Khóa ngoại tham chiếu bảng User \\ \hline
\code{destination} & String & Có & Tỉnh/Thành phố của chuyến đi \\ \hline
\code{duration} & Number & Có & Số ngày hành trình (1--7 ngày) \\ \hline
\code{total\_budget} & Number & Có & Tổng ngân sách dự kiến (VNĐ) \\ \hline
\code{people} & Number & Có & Số lượng thành viên tham gia \\ \hline
\code{days} & Array[Object] & Có & Danh sách các ngày, mỗi ngày gồm các hoạt động \\ \hline
\code{createdAt} & Date & Có & Thời gian tạo lập chuyến đi \\ \hline
\end{tabular}
\end{table}

\section{Giải thuật Visited Set Graph (Kiểm soát Không Trùng Lặp)}

\begin{lstlisting}[language=JavaScript, caption=Giải thuật kiểm soát không lặp địa điểm (aiService.js)]
// Algorithm: Visited Set Graph for Non-Repeating Itinerary
function taoLichTrinhThongMinh(destination, duration, budget, people, selectedPlaces, dbPlaces) {
    const usedPlaceNames = new Set();
    
    // Step 1: Mark user-selected places into priority queue
    selectedPlaces.forEach(name => usedPlaceNames.add(name));
    
    function layDiaDiemKhongTrung(availableList, fallbackName, type, cost) {
        // Filter out any place that has already been visited
        const unvisited = availableList.filter(p => !usedPlaceNames.has(p.name));
        if (unvisited.length > 0) {
            const picked = unvisited[0];
            usedPlaceNames.add(picked.name);
            return picked;
        }
        // Fallback generator with unique incremental ID
        const fallback = { 
            name: `${fallbackName} (Điểm ${usedPlaceNames.size + 1})`, 
            type: type, 
            cost: cost 
        };
        usedPlaceNames.add(fallback.name);
        return fallback;
    }
    
    // Step 2: Allocate slots for each day (Breakfast, Morning Attraction, Lunch, Check-in, Cafe, Dinner)
    const itinerary = [];
    for (let day = 1; day <= duration; day++) {
        itinerary.push(buildDailySchedule(day, layDiaDiemKhongTrung, dbPlaces));
    }
    return itinerary;
}
\end{lstlisting}
